\noindent\hspace*{1.2em}\kbvlabel{ಉದಾ}\hspace{0.6em}ನಿರವಧಿಕ ನಿರ್ಮಾಯ ನಿಸ್ಪೃಹ ನಿರುಪಮಿತ ನಿರ್ದ್ವಂದ್ವ ನಿರ್ಗುಣ\par

\begin{kbvverse}
\noindent ನಿರವಯವ ನಿರ್ಲೇಪ ನಿರಾಗ್ರಹ ನಿರಾಧಾರ ।
\end{kbvverse}

\indent ನಿರುಪಮ ನಿರಾಮಯ ನಿರಂತರ ನಿರವಶೇಷ ನಿರಂಗ ನಿರ್ಮಲ\par

\begin{kbvverse}
\noindent ನಿರತಿಶಯ ನಿಷ್ಕಳ ಮಹೇಶ್ವರ ಕರುಣಿಸುವುದೆಂದ ॥
\end{kbvverse}

\indent (ಕರ್ಣಾಟ ಭಾರತ ಕಥಾಮಂಜರಿ : ಅರಣ್ಯಪರ್ವ) ಇದು ವೃತ್ತಾನುಪ್ರಾಸಕ್ಕೆ ಸೊಗಸಾದ ಉದಾಹರಣೆ. ಇಲ್ಲಿ ನಕಾರವು ಹಲವಾರು ಬಾರಿ ಪುನರುಕ್ತವಾಗಿರುವುದರಿಂದ ಪದ್ಯದ ಮಾಧುರ್ಯ ಹೆಚ್ಚಿದೆ.\par

\kbvsechead{(೨)\hspace{0.45em}ಛೇಕಾನುಪ್ರಾಸ:}

\begin{kbvverse}
\noindent ರುಣ ನಿರೋಧಿ ವಿಭಾಡ ರುಣ ರುಣ । \\
\noindent ರುಣು ಜಗತ್ರಯ ಜಟ್ಟಿ ರುಣ ರುಣ । \\
\noindent ರುಣು ರುಣೆಂಬಬ್ಬರಣೆ ಮಸಗಿದುದೆರಡು ಬಾಹೆಯಲಿ ॥
\end{kbvverse}

\indent ಹೀಗೆ, ಎರಡೆರಡು ವ್ಯಂಜನಗಳಿಂದ ಕೂಡಿದ ಶಬ್ದಗಳ ಪುನರಾವರ್ತನೆ ಯಿಂದ ಪದ್ಯದ ಮಾಧುರ್ಯ ಹಾಗು ಸೌಂದರ್ಯ ಹೆಚ್ಚಿದಲ್ಲಿ ಅದು ಛೇಕಾನು ಪ್ರಾಸವೆನಿಸುವುದು.\par

\vspace{3pt plus 2pt}\noindent\hspace*{1.2em}\kbvlabel{ಉದಾ}\par\vspace{1pt}

\begin{kbvverse}
\noindent ರಾಮನಂ ಭುವನಾಭಿ ರಾಮನಂ ಗುಣರತ್ನ । \\
\noindent ಧಾಮನಂ ರಘುಕುಲೋದ್ಧಾಮನಂ ರೂಪಜಿತ । \\
\noindent ಕಾಮನಂ ಸತ್ಕೀರ್ತಿ ಕಾಮನಂ ಶರಣಜನವಾರ್ಧಿಯಂ ಮಿಗೆ ಪೆರ್ಚಿನ । \\
\noindent ಸೋಮನಂ ಸೌಭಾಗ್ಯ ಸೋಮನಂ ಕುವಲಯ । \\
\noindent ಶ್ಯಾಮನಂ ನಿಜತನು ಶ್ಯಾಮನಂ ಘನಪುಣ್ಯ । \\
\noindent ನಾಮನಂ ಸಂತತಂ ನಾ ಮನಂದಣಿಯೆ ರಮಿಸದೆ ಬಾಳೆನೆಂದಳು ।
\end{kbvverse}

\noindent\hfill{\fontsize{9.5pt}{12.5pt}\selectfont\color{kbvgrey}--ಕನ್ನಡ ಜೈಮಿನಿ ಭಾರತ : ಸೀತಾ ಪರಿತ್ಯಾಗ ವೃತ್ತಾಂತ}\par\vspace{3pt plus 2pt}

\indent ಇಲ್ಲಿ "ಮನಂ" ಎಂಬುದು ಹಲವು ಬಾರಿ ಪುನರುಕ್ತವಾಗಿರುವ ದೆಸೆಯಿಂದ ಪದ್ಯದ ಮಾಧುರ್ಯ ಹಾಗೂ ಸೌಂದರ್ಯ ಹೆಚ್ಚಿದೆ ; ಇದು ಛೇಕಾನು ಪ್ರಾಸಕ್ಕೆ ಉತ್ತಮ ಉದಾಹರಣೆ.\par
